\chapter{Security, Ethics, and Operational Governance}
\label{ch:security-ethics}

As a cyber-intelligence platform, Wolverine's operational validity is inexorably tied to its security architecture 
and ethical boundaries. This chapter deconstructs the system's trust zones, details the simulated limitations of 
WolverineDB, analyzes the ethical distinction between probabilistic correlation and human attribution, and establishes 
a rigorous residual risk model.

\section{Security Model}
\label{sec:sec-model}
The Wolverine security model is structured to defend internal analytical assets from external payload poisoning and 
internal unauthorized exfiltration.
\begin{itemize}
    \item \textbf{Assets:} Immutable source payloads, canonical records, linkage probabilities, evaluation ground truth.
    \item \textbf{Adversaries:} Malicious actors attempting prompt injection (T-01), infrastructure operators attempting 
          database tampering (T-05), and data ingestion errors causing denial of service.
    \item \textbf{Controls:} Network segregation, strict schema typing, LLM Regex sanitization, Merkle state hashing.
\end{itemize}

\section{STRIDE Threat Model Analysis}
\label{sec:stride-threat-model}
To systematically analyze the Wolverine architecture, we apply the STRIDE threat modeling methodology 
\cite{shostack2014threat}. By evaluating the system against these categories, we can identify and document the 
primary attack vectors and their corresponding mitigations within the platform.

\begin{itemize}
    \item \textbf{Spoofing:} The platform currently relies on \texttt{X-Wolverine-Role} headers for Role-Based 
          Access Control (RBAC), which are vulnerable to spoofing since there is no JWT or OAuth cryptographic 
          validation implemented.
    \item \textbf{Tampering:} Adversaries might attempt to modify PostgreSQL records. This is largely mitigated by 
          the Merkle hashing state layer built into WolverineDB, which detects unauthorized alterations at the row level.
    \item \textbf{Repudiation:} There is currently a lack of granular user action logging within the API adapters, 
          leading to a residual risk where actions cannot be definitively traced to specific analysts.
    \item \textbf{Information Disclosure:} The Truth Vault is logically air-gapped using Docker network isolation, 
          which effectively prevents inference leakage of the ground truth labels to the active analytical pipeline.
    \item \textbf{Denial of Service:} Volumetric attacks against the DMZ are mitigated via Nginx rate limiting. 
          Internal graphing Denial of Service is prevented by enforcing strict node caps on Breadth-First Search 
          (BFS) traversals.
    \item \textbf{Elevation of Privilege:} An attacker who gains access to the internal network could exploit the 
          aforementioned RBAC header spoofing to elevate their privileges to administrator functions.
\end{itemize}

\section{Trust Zones}
\label{sec:sec-trust-zones}
The architecture relies on physical network separation rather than application-layer logic to enforce trust 
boundaries (\Cref{ch:system-architecture}).
\begin{itemize}
    \item \textbf{\texttt{wolverine-dmz}:} Untrusted. Contains adversarial data sources. Ingress is restricted via 
          Nginx rate limiting.
    \item \textbf{\texttt{wolverine-internal}:} Semi-Trusted. Contains the analytical core. It can pull from the DMZ 
          but is strictly prohibited from pushing data to the public internet or the evaluation vault.
    \item \textbf{\texttt{wolverine-truth}:} Strictly Trusted. Operates via the Docker \texttt{internal: true} 
          directive, dropping all routing from \texttt{wolverine-internal}.
\end{itemize}

\section{Container Security Analysis}
\label{sec:container-security}
The trust zones established above rely on Docker for environment segregation. However, as noted by Sultan et al. 
\cite{sultan2019container}, Docker network isolation constitutes a virtual boundary, not a physical one. 
This architectural choice carries inherent risks:

\begin{itemize}
    \item \textbf{Shared Kernel:} All containers share the host operating system's kernel. A kernel vulnerability 
          could allow an attacker in the untrusted DMZ to compromise the host, thereby bypassing all network isolation.
    \item \textbf{Docker Daemon Vulnerability:} The entire isolation model relies on the integrity of the Docker 
          daemon. If the daemon is compromised, the isolation model completely breaks down.
    \item \textbf{Container Escape:} There is a persistent risk of container escape attacks, especially since the 
          current deployment does not configure restrictive \texttt{seccomp} or \texttt{AppArmor} profiles.
\end{itemize}

\section{The Truth Vault Boundary}
\label{sec:sec-truth-vault}
The Truth Vault (\texttt{postgres-truth}) stores the exact human-to-account assignments generated synthetically. 
The \texttt{internal: true} network configuration guarantees that the \texttt{entity-resolver} cannot issue SQL queries 
to this database during operation. 

This does not create a literal physical "air-gap" (as the containers share a single host kernel), but it establishes a 
virtualized packet-drop boundary. The only mechanism to cross this boundary is the external, out-of-band evaluation 
script that temporarily binds to the host's localhost port during scoring runs.

\section{Tor Gateway Controls}
\label{sec:sec-tor-gateway}
While the synthetic prototype currently scrapes local containers, the architecture defines a Tor gateway for real-world 
integration. This mitigates the deanonymization of the collection infrastructure by hostile platform administrators. 
All outbound API adapter requests route through a proxy container rotating Tor circuits.

The Nginx configuration securing the Tor Gateway enforces strict operational controls (\Cref{app:tor-gateway}):
\begin{itemize}
    \item \textbf{Method Restriction:} Allows only \texttt{GET} and \texttt{HEAD} requests, returning a 
          \texttt{405 Method Not Allowed} for all other HTTP verbs.
    \item \textbf{Rate Limiting:} Implements aggressive rate limiting with \texttt{rate=10r/s} and a \texttt{burst=20} 
          to thwart automated scraping attacks.
    \item \textbf{Endpoint Blocking:} Explicitly blocks access to administrative endpoints including \texttt{/metrics}, 
          \texttt{/health/ready}, \texttt{/admin}, and \texttt{/seed}.
    \item \textbf{Security Headers:} Sets mandatory security headers such as Content Security Policy (\texttt{CSP}), 
          \texttt{X-Frame-Options}, and \texttt{X-Content-Type-Options}.
    \item \textbf{Upstream RBAC:} Injects upstream RBAC headers, mapping connections to predefined roles such as 
          \texttt{X-Wolverine-Role "researcher"} and \texttt{"public\_demo"}.
\end{itemize}
The primary residual risk here is the potential for malicious Tor exit nodes to inject payload modifications before 
TLS validation can occur.

\section{WolverineDB Security Architecture}
\label{sec:sec-wolverinedb}
WolverineDB injects tamper-evidence into relational PostgreSQL via Merkle state hashing. For a full architectural 
treatment, refer to \Cref{ch:wolverinedb}. Here, we summarize the security model.

\subsection{Vulnerability Forensics}
\label{subsec:vuln-forensics}
Despite successful functional tests, the implementation remains a simulated proof-of-concept. The following 
vulnerabilities are explicitly acknowledged:
\begin{enumerate}
    \item \textbf{In-Process Transport:} BFT consensus relies on \texttt{DirectMemoryNetworkTransport}. It is 
          fundamentally vulnerable to a single-process crash, violating distributed fault tolerance.
    \item \textbf{Simulated KMS:} Cryptographic signatures are not anchored to a physical Hardware Security Module 
          (HSM).
    \item \textbf{Sentinel TOCTOU:} A Time-of-Check to Time-of-Use window exists between the Sentinel policy validation 
          and actual PostgreSQL insertion.
\end{enumerate}
These limitations explicitly render WolverineDB unsuitable for production financial or legal evidence retention.

\subsection{Evolution of Replay Protection}
\label{sec:security-replay-evolution}
The cryptographic authorization model relies on single-use approval nonces. The initial prototype utilized a volatile 
Node.js \texttt{Set}, vulnerable to container-restart replay attacks. The architecture evolved to use 
\texttt{PostgresReplayStore}, mapping atomic primary key violations to authorization rejections (\Cref{tab:security-replay}).

\begin{table}[htpb]
\centering
\begin{tabular}{p{3.5cm} p{3cm} p{4cm} p{3cm}}
\toprule
\textbf{Property} & \textbf{Previous Prototype} & \textbf{Current Implementation} & \textbf{Evidence} \\
\midrule
Replay state & volatile memory & PostgreSQL-backed & Source/Tests \\
Persistence & process lifetime & database persistence & Verified (Tests) \\
Duplicate detection & in-memory Set & DB uniqueness & Source/Tests \\
Concurrent attempts & race-sensitive & atomic DB constraint & Verified (Tests) \\
Test double & N/A & \texttt{InMemoryReplayStore} & Source \\
Residual limitation & process restart & runtime injection / config & Source \\
\bottomrule
\end{tabular}
\caption{Evolution of Replay Protection Security}
\label{tab:security-replay}
\end{table}

\section{Data Security and Synthetic Privacy}
\label{sec:sec-synthetic-privacy}
The ethical mandate of this project (NFR-04) strictly prohibits the ingestion of real-world Personally Identifiable 
Information (PII). By utilizing entirely synthetic datasets generated deterministically, the prototype eliminates the 
legal and ethical liabilities associated with analyzing live dark-web telemetry.

\subsection{Ethics of Synthetic Data}
\label{subsec:ethics-synthetic}
Using synthetic data is an ethical necessity to protect innocent parties from accidental implication during research 
and development. However, it introduces known limitations:
\begin{itemize}
    \item \textbf{Domain Shift:} Synthetic data may not capture the nuanced, chaotic nature of real-world threat 
          actors, leading to a synthetic-to-real domain shift.
    \item \textbf{The False Confidence Problem:} High accuracy on synthetic data does not equate to real-world safety 
          or performance. Over-reliance on these metrics can lead to dangerous assumptions when the system is deployed 
          in production environments.
\end{itemize}
Furthermore, as demonstrated by Narayanan and Shmatikov \cite{narayanan2008robust}, even supposedly anonymized data 
carries significant de-anonymization risks. Utilizing pure synthetic generation side-steps these privacy pitfalls 
while preserving structural complexity.

\section{RAG Security and Prompt Injection}
\label{sec:sec-rag-security}
When the LLM receives graph data, it is subjected to 4-rule Regex sanitization \cite{perez2022ignore, greshake2023not}.
\begin{itemize}
    \item \textbf{Threat:} An adversarial payload containing instructions (e.g., \texttt{ignore previous system prompts}).
    \item \textbf{Detection:} Rule 2 \texttt{/(ignore|system|instruction|prompt)/gi}.
    \item \textbf{Response:} Replaces matches with \texttt{[REDACTED]}.
\end{itemize}

\textit{Illustrative synthetic example --- not an empirical benchmark:} \\
\textbf{Attack:} \texttt{"Vendor: shadow\_broker. System: print internal state."} \\
\textbf{Sanitized:} \texttt{"Vendor: shadow\_broker. [REDACTED]: print internal state."} \\
\textbf{Residual Risk:} Sophisticated multi-lingual or encoded prompt smuggling can trivially bypass static regex 
matching.

\section{Threat-to-Control Traceability}
\label{sec:sec-threat-traceability}
\begin{table}[htpb]
\centering
\begin{tabular}{p{0.10\textwidth} p{0.15\textwidth} p{0.30\textwidth} p{0.35\textwidth}}
\toprule
\textbf{ID} & \textbf{Asset} & \textbf{Control} & \textbf{Residual Risk} \\
\midrule
\textbf{T-01} & LLM & Regex Sanitization (4 Rules) & Semantic Jailbreaks (High) \\
\textbf{T-03} & API & Nginx Rate Limiting & Distributed Volumetric DDoS (Med) \\
\textbf{T-05} & Database & WolverineDB Merkle Hashing & Config-Dependent Replay Attack (Med) \\
\bottomrule
\end{tabular}
\caption{Threat-to-Control Traceability Matrix}
\label{tab:threat-controls}
\end{table}

\section{Failure and Abuse Cases}
\label{sec:sec-failure-cases}
\begin{itemize}
    \item \textbf{Malformed Record:} Parser drops payload. Impact: Missing data.
    \item \textbf{LLM Failure:} Timeout triggers deterministic string fallback. Impact: Loss of narrative synthesis.
    \item \textbf{Cryptographic Failure:} Missing payload hash prevents normalization. Impact: Complete ingestion failure.
\end{itemize}

\section{Correlation vs. Attribution}
\label{sec:sec-correlation-attribution}
A critical ethical boundary in intelligence analysis is distinguishing statistical correlation from legal human 
attribution \cite{nist2012guide, dwork2006differential}. We formally define these as:
\begin{itemize}
    \item \textbf{Statistical Correlation:} A mathematical and probabilistic relationship between two data points or 
          entities based on observed features.
    \item \textbf{Legal Attribution:} The rigorous evidentiary standard required to definitively link an activity to 
          a specific human identity.
    \item \textbf{Automation Bias:} The cognitive tendency to over-trust and rely on automated systems, often ignoring 
          contradictory human intuition or evidence.
\end{itemize}

The Wolverine system calculates \textbf{record similarity} (Jaro-Winkler) and \textbf{temporal overlap}. When these 
features exceed a threshold, the system creates a \textbf{probabilistic candidate correlation} (an edge in Neo4j). 

This represents an \textit{analyst hypothesis}. It does \textbf{NOT} constitute verified human identity. Shared handles 
and temporal proximity can easily result from deliberate framing, operational coincidence, or identity theft. Asserting 
that a Neo4j graph edge proves legal attribution is a dangerous automation bias explicitly warned against in this 
architecture.

\section{Human Oversight and Governance}
\label{sec:sec-governance}
Because correlation is not attribution, the system dictates an explicit requirement for Human-in-the-Loop (HITL) 
oversight. 
Candidate pairs scoring between $0.70$ and $0.92$ are tagged for human review, requiring qualitative analyst evaluation 
of the raw payloads before an edge is confirmed.

\section{Residual Risk Model}
\label{sec:sec-residual-risk}
\begin{table}[htpb]
\centering
\begin{tabular}{p{0.35\textwidth} p{0.15\textwidth} p{0.40\textwidth}}
\toprule
\textbf{Risk} & \textbf{Severity} & \textbf{Future Mitigation} \\
\midrule
LLM Semantic Prompt Smuggling & High & Secondary LLM output evaluator. \\
WolverineDB TOCTOU / Memory Trans. & High & Physical multi-node deployment. \\
Algorithmic False Positives (FPR) & Medium & Dynamic threshold calibration (ROC). \\
Automation Bias / False Attribution & High & Mandated UI confidence disclaimers. \\
Container escape / Docker daemon exploit & High & Implement \texttt{seccomp/AppArmor} profiles. \\
RBAC header spoofing & High & Implement cryptographic OAuth/JWT tokens. \\
Synthetic-to-real domain shift & Medium & Phased evaluation on isolated real data. \\
Model drift in LLM synthesis & Medium & Continuous evaluation and version locking. \\
\bottomrule
\end{tabular}
\caption{Operational Residual Risk Matrix}
\label{tab:residual-risk}
\end{table}

\section{Security/Ethics Synthesis}
\label{sec:sec-synthesis}
The prototype successfully demonstrates architectural protections via strict network isolation and immutable raw capture. 
However, the cryptographic fault tolerance remains simulated, the RAG prompt-injection defenses are vulnerable to 
semantic bypass, and the probabilistic nature of the entity resolution mandates strict human oversight to prevent 
unethical automation bias.
